% ---------------------------------------------------------------
%  Section IX -- Limitations, Ethics, and Privacy
% ---------------------------------------------------------------

\subsection{Data Access and Operator Dependencies}

The most significant practical limitation is the absence of machine-readable
transit feeds for Baku Metro and BakuBus.
Both operators would need to provide GTFS static schedules and, ideally,
GTFS-RT trip updates under a formal data-sharing agreement before IRIDIUM
can offer production-quality multimodal journey planning.
Traffic data require a commercial license and an API key; without these,
dynamic edge state $\mathbf{d}_e(t)$ is absent and routing falls back to
free-flow estimates via Eq.~(\ref{eq:heuristic_forecast}).
Event and energy pipelines are documented but lack implemented adapters;
both appear in Table~\ref{tab:providers} with \texttt{not\_impl.}\ status.
Consent-based GNSS telemetry depends on Traccar deployment, device enrollment,
and informed consent from participants.
As a result, the current evaluation can confirm only OSM network ingestion
and Open-Meteo weather integration; claims about forecasting accuracy,
multimodal routing quality, equity scores, and anomaly detection performance
remain deferred pending resolution of these dependencies.

\subsection{Evaluation Scope and Generalizability}

The absence of trained predictive models means that no comparison against
published ST-GNN benchmarks (DCRNN, ST-GCN, Graph WaveNet, ASTGCN)
can be provided at this stage.
Response-latency measurements (Section~\ref{sec:results}) were obtained
on a loopback interface under zero concurrent load and represent a lower
bound on production latency; production deployments will experience
additional overhead from network round-trips, PostgreSQL replication,
connection pooling, and concurrent request handling.
All hyperparameters in Table~\ref{tab:hyperparams} are design targets
rather than experimentally validated values.

\subsection{Mobility Equity Score: Bias and Scope Limitations}

The $\MES_d$ (Eq.~\ref{eq:mes}) is sensitive to indicator selection,
district boundary definitions, and spatial data coverage.
Under-coverage of informal settlements or peripheral areas can bias
normalized indicators $z_{d,m}$ downward for those districts, artificially
compressing $\Delta_{\MES}$ and $\mathrm{Var}(\MES)$.
Within-district inequality is invisible by construction: a high $\MES_d$
for a district with a spatially concentrated transport hub does not imply
equitable access for all residents.
The score is a derived analytic index produced by the platform analytics
layer; it is not validated against an independent ground truth, and
policy decisions must account for these limitations and seek external
verification~\cite{karner2025transportation_equity}.

\subsection{Privacy, Consent, and Telemetry}

Consent-based GNSS telemetry involves personal location data subject to
national data protection legislation in Azerbaijan and principles aligned
with the General Data Protection Regulation (GDPR) for any processing
involving EU-resident individuals.
The IRIDIUM design keeps raw traces out of the central server through
short-lived in-memory processing and federated model training; no personal
GNSS traces are persisted in the central repository.
Operators deploying the telemetry layer must conduct a data-protection
impact assessment (DPIA), define retention limits consistent with the
principle of data minimization, implement a consent withdrawal mechanism,
and provide a documented right-to-erasure procedure.

\subsection{Differential Privacy in Federated Learning}

The FL pathway uses FedAvg aggregation (Eq.~\ref{eq:fedavg}).
Without differential privacy, gradient inversion attacks could potentially
reconstruct training samples from aggregated parameter
updates~\cite{kairouz2021advances}.
The planned noise calibration via Eq.~(\ref{eq:dp_noise}) requires
specifying the privacy budget $(\varepsilon,\delta)$ before deployment.
The privacy-utility trade-off (i.e., the degradation in forecasting
accuracy $\Delta\mathrm{MAE}$ as a function of $\varepsilon$) has not yet
been characterized, and must be evaluated empirically for the Azerbaijani
traffic data distribution.
Deployers must conduct this analysis and assess the required privacy budget
before enabling FL in production.

\subsection{Model Generalization and Transferability}

All design choices, default hyperparameters, and provider configurations
are calibrated for Azerbaijani cities and the data sources described in
Section~\ref{sec:data}.
Applying the platform to other countries requires re-running the OSM
import, registering locally relevant GTFS and traffic providers, and
potentially retraining the ST-GNN on city-specific historical data.
No formal transferability study comparing forecast accuracy across
multiple cities or road-network topologies has been conducted.

\subsection{Deployment Risks and Operational Safety}

Anomaly detection based on z-score thresholds will produce false positives
under heavy distributional shift (e.g., a public holiday reducing baseline
traffic) and false negatives when incidents develop gradually below the
threshold $\tau_S$.
The false-positive rate estimate from Eq.~(\ref{eq:fpr}) assumes approximate
normality of the rolling speed distribution and independence across time
steps; neither assumption holds perfectly under real traffic dynamics.
Real-time routing decisions may briefly direct users toward congested
segments if an incident is detected with a lag of up to $w$ time steps
(default: 12 steps, equivalent to one hour at five-minute polling).
The system is best-effort and must not serve as the sole basis for
safety-critical operational decisions.
Incident escalation procedures and operator fallback protocols are
documented in the operational guide accompanying the repository.
